\section{Short-window freezing and actual phase packets}\label{sec:packets}

This section develops the result of BB Note25~\cite{BB25}, accepted
without mathematical correction in the E11 review~\cite{BBReview11}.

Fix exactly
\begin{equation}\label{eq:short-window}
 \begin{gathered}
 h_U=U^{1/32},\qquad Y\in[U+h_U,2U-h_U],\\
 \dd\nu_{h_U}(s)=h_U^{-1}(1-|s|/h_U)_+\dd s,\qquad
 A_U(Y)=\int Q(Y+s)\dd\nu_{h_U}(s).
 \end{gathered}
\end{equation}
The center measure has mass one. The fixed question is whether
$A_U(Y)\ge5/66^2$ at admissible actual centers along an unbounded
sequence of $U$. Positivity would then force a violation of the
every-center $4/66^2$ budget somewhere in the window.
The results of this section do not decide that question.

\subsection{Freezing in the consumed Hilbert norm}

All norms in this subsection use
$\mathcal H_U=L^2(\dd\nu_{h_U}(s)P(v)\dd v)$.
At the fixed center $Y$, write $I_Y=I_Y^{\rm m}$ and set
\[
 G_Y(s,v)=\sum_{n\in I_Y}w_n(Y)
      \ee^{\ii(\tau_Y+(s+v)/2)\log n},\qquad
 F_U(Y)=\|G_Y\|_{\mathcal H_U}^2.
\]
\begin{theorem}[Uniform square-root transfer]\label{thm:freezing}
For every admissible $Y$, uniformly in the fractional parts of both
real cutoff values,
\begin{equation}\label{eq:freezing}
 \|M_{Y+s}-G_Y\|_{\mathcal H_U}\ll U^{-771/16384},\qquad
 |\sqrt{A_U(Y)}-\sqrt{F_U(Y)}|\ll U^{-1/8192}.
\end{equation}
The frozen form has the exact expansion
\begin{equation}\label{eq:frozen-form}
 F_U(Y)=\sum_{n,m\in I_Y}w_n(Y)w_m(Y)\ee^{\ii\tau_Y\log(n/m)}
 \widehat P\left(\frac{\log(n/m)}2\right)
 \sinc^2\left(\frac{h_U\log(n/m)}4\right),
\end{equation}
where $\sinc q=\sin q/q$ with value one at zero.
\end{theorem}

\begin{proof}
For either cutoff exponent $p\in\{3/64,5/16\}$, the real function
$\ee^{pL_X}$ has derivative $p\ee^{pL_X}/X$. Its excursion across
the whole window is $O(U^{p-31/32})=o(1)$.
Since the function is monotone, each floor can change by at most one
integer. At a changed endpoint, $\log n=pL_Y+O(1)$ and
$L_X-L_Y=O(h_U/U)$, so
\[
 w_n(X),w_n(Y)\ll U^{-65p/64+p^2/4}.
\]
The two endpoint powers are $-771/16384$ and $-75/256$.
This also bounds terms that occur only in the frozen sum.

For common indices,
\[
 \partial_X\log w_n(X)=O(U^{-1}),\qquad
 \tau_{Y+s}-\tau_Y-s/2
 =\frac\pi8(L_{Y+s}^{-1}-L_Y^{-1})
 =O\left(\frac{h_U}{U\log^2U}\right).
\]
After multiplication by $\log n$, the phase error is
$O(h_U/(U\log U))$. On the middle band,
$w_n(Y)\le n^{-15/16}$ and $\sum_{I_Y}w_n(Y)\ll U^{5/256}$.
The total common-term error is therefore
$O((h_U/U)U^{5/256})=O(U^{-243/256})$, uniformly for all real $v$.
Adding the two endpoint terms proves the first norm bound.
Theorem~\ref{thm:localization}, integrated against the probability
measure $\nu_{h_U}$, pays the raw-to-middle norm by $U^{-1/8192}$.
The reverse triangle inequality proves the second bound.

The normalized indicator of $[-h_U/2,h_U/2]$ has Fourier transform
$\sinc(h_U\omega/2)$; its self-convolution has precisely the triangular
density in~\eqref{eq:short-window}. In the finite squared sum,
the $s$ and $v$ frequencies are both $\log(n/m)/2$.
Their transforms give~\eqref{eq:frozen-form}. Every ordered pair is
retained. The form is nonnegative because it is a squared joint norm,
not because its signed cross terms have a fixed sign.
\end{proof}

The transfer is deliberately between square-root moments. It is not
an additive $o(1)$ estimate for the moments themselves. The full norm
bound proved next is $O(U^{1/256})$; multiplying it by the localization
error would give $U^{31/8192}$, which grows. A fixed strict surplus
in the square-root norm can be transferred; an equality-level bound
requires its error payment.

\subsection{Summable logarithmic cells and sparse deletion}

Partition the middle indices into cells
\[
 \mathcal C_j=\{n\in I_Y:j\le h_U\log n<j+1\},\qquad
 a_j=\sum_{n\in\mathcal C_j}w_n(Y).
\]
For $n,m$ in cells $j,k$, the absolute multiplier in
\eqref{eq:frozen-form} is at most $C(1+|j-k|)^{-2}$.
Indeed for $|j-k|\ge2$ one has
$h_U|\log(n/m)|\ge|j-k|-1$ and
$|\sinc q|\le\min(1,|q|^{-1})$; the bound
$|\widehat P|\le2/5$ covers the remaining cases.
Since these rows are summable,
\begin{equation}\label{eq:cell-row}
 F_U(Y)\le C\sum_ja_j^2.
\end{equation}
This follows by summing $2a_ja_k\le a_j^2+a_k^2$.
It applies to any subset of indices and even to arbitrary fixed
coefficient arguments, so it is an upper comparison only.

Put $D_j=\ee^{j/h_U}$. The number of integers in a cell is at most
$D_j(\ee^{1/h_U}-1)+1$. In every occupied cell,
$D_j/h_U\gg U^{1/64}\to\infty$, because the lower middle scale is
$U^{3/64}$. Thus the $+1$ is absorbed and the number is $O(D_j/h_U)$.
Using $w_n\le n^{-15/16}$ gives
\begin{equation}\label{eq:cell-mass}
 a_j\ll D_j^{1/16}/h_U,\qquad
 \max_ja_j^2\ll h_U^{-2}U^{5/128}=U^{-3/128}.
\end{equation}
If $j_*$ is the last occupied cell, the full geometric sum satisfies
\[
 \sum_{j\le j_*}\ee^{j/(8h_U)}
 =\frac{\ee^{j_*/(8h_U)}}{1-\ee^{-1/(8h_U)}}
 \ll h_U U^{5/128}.
\]
Consequently~\eqref{eq:cell-row} and~\eqref{eq:freezing} yield
\begin{equation}\label{eq:all-cell-bound}
 F_U(Y)\ll U^{1/128},\qquad A_U(Y)\ll U^{1/128}.
\end{equation}
The triangular window removes the logarithm from the earlier
short-center allowance, but the new allowance still grows.

\begin{theorem}[Sparse-packet effect on the complete moment]\label{thm:sparse}
Let $S\subset I_Y$ meet at most $r\ge1$ cells.
Its frozen standalone energy is $O(rU^{-3/128})$.
Let $A_U^{\setminus S}(Y)$ denote the actual average after deleting
these fixed indices from the complete raw sum at every center in the
window. Then
\begin{equation}\label{eq:sparse-full}
 |A_U(Y)-A_U^{\setminus S}(Y)|
 \ll\sqrt r\,U^{-1/128}+rU^{-3/128}.
\end{equation}
In particular its complete effect tends to zero for $r=o(U^{1/64})$.
\end{theorem}

\begin{proof}
The standalone frozen bound follows from~\eqref{eq:cell-row} and
\eqref{eq:cell-mass}. The whole fixed set $I_Y$ lies strictly inside
the raw cutoffs throughout the window eventually, because
$1/32<3/64<5/16<1/2$ and $L_{Y+s}/L_Y\to1$ uniformly.
There is no moving membership for the selected packet
\[
 P_{S,X}(v)=\sum_{n\in S}w_n(X)\ee^{\ii(\tau_X+v/2)\log n}.
\]
The common-term freezing proof pays its error by $O(U^{-243/256})$.
Its actual joint norm is therefore $O(\sqrt r\,U^{-3/256})$.
For $V=T_{Y+s}$ and $W=P_{S,Y+s}$ in $\mathcal H_U$,
\[
 \|V\|^2-\|V-W\|^2=2\Re\ip VW-\|W\|^2.
\]
Using the complete bound $\|V\|\ll U^{1/256}$ from
\eqref{eq:all-cell-bound} proves~\eqref{eq:sparse-full}.
This includes the interaction with every low, middle and high index.
\end{proof}

The standalone energy vanishes in the different, larger sufficient
range $r=o(U^{3/128})$. Conversely, a fixed positive standalone energy
would require $r\gg U^{3/128}$ under this estimate.
These statements must not be interchanged with the complete-perturbation
range. None bounds the energy of the untouched remainder below the
witness threshold.

\subsection{Quartic integer alignment at actual heights}

We now test one concrete lower mechanism without replacing the physical
phases by freely chosen arguments. Let $D\to\infty$ through integers,
and choose the unique sufficiently large real $Y_D$ satisfying
\begin{equation}\label{eq:quartic-center}
 \tau_{Y_D}=24\pi D^4.
\end{equation}
Existence and uniqueness on the large branch follow from continuity,
$\tau_X'\ge1/4$ and $\tau_X\to\infty$.
Then $Y_D\sim48\pi D^4$. Set $U=2Y_D/3$, so that the center is
admissible, and keep $h_U=U^{1/32}\asymp D^{1/8}$.
Take the fixed packet
\[
 J=\lfloor D^{1/5}/4\rfloor,\qquad S_D=\{D,D+1,\ldots,D+J\}.
\]
Uniformly across the window, $L_X=4\log D+\log12+o(1)$ and
$\log(D+j)/L_X\to1/4$. Thus the packet belongs to both middle and
raw populations throughout the window, with room to spare at both
integer cutoffs.

\begin{theorem}[An aligned packet with vanishing complete effect]\label{thm:quartic}
For the actual centers~\eqref{eq:quartic-center},
\begin{equation}\label{eq:quartic-energy}
 \|P_{S_D,Y_D+s}\|_{\mathcal H_U}^2
 \asymp D^{-241/160}\asymp U^{-241/640},
\end{equation}
and
\begin{equation}\label{eq:quartic-full}
 |A_U(Y_D)-A_U^{\setminus S_D}(Y_D)|\ll U^{-59/320}.
\end{equation}
\end{theorem}

\begin{proof}
Integrating the geometric-series remainder for $1/(1+x)$ gives
\begin{equation}\label{eq:quartic-taylor}
 \begin{split}
 24\pi D^4\log(1+j/D)
 &=2\pi(12D^3j-6D^2j^2+4Dj^3-3j^4)+R_j,\\
 0\le R_j&\le\frac{24\pi}{5}\frac{j^5}{D}
                  \le\frac{3\pi}{640}<\frac1{32}.
 \end{split}
\end{equation}
The polynomial in parentheses is an integer.
For $|X-Y_D|\le h_U$ and $|v|\le1$, the additional relative phase
against the term $n=D$ is
\[
 (\tau_X-\tau_{Y_D}+v/2)\log(1+j/D)
 =O((h_U+1)J/D)=O(D^{-27/40}).
\]
Eventually all packet terms therefore have relative arguments within
$1/16$ modulo $2\pi$ on this entire portion of the joint domain.
Their weights satisfy $w_{D+j}(X)\asymp D^{-61/64}$, since
$(\log(D+j))^2/(4L_X)=\log D/16+O(1)$ uniformly.
After rotation by the first term's phase, real projection gives a
lower bound $\cos(1/16)\sum_{S_D}w_n(X)$ on $|v|\le1$.
The triangle inequality gives the corresponding upper bound for all
$v$. The restricted Poisson mass is positive and the center measure
has mass one. Hence the energy has order
\[
 (D^{1/5}D^{-61/64})^2=D^{2/5-61/32}=D^{-241/160},
\]
proving~\eqref{eq:quartic-energy}.
Finally apply the exact squared-norm difference used in
Theorem~\ref{thm:sparse}, now with the sharper packet norm:
\[
 |A_U-A_U^{\setminus S_D}|
 \ll U^{1/256}U^{-241/1280}+U^{-241/640}
 \ll U^{-59/320}.
\]
The first factor is the norm of the complete raw sum, so the whole
background interaction is included.
\end{proof}

The packet also meets at most two logarithmic cells eventually,
because $h_U\log(1+J/D)=o(1)$.
It succeeds at actual alignment and fails at fixed mass. This stops
this isolated quartic-packet mechanism and the stated sparse deletion
class. It neither disproves the short-center witness nor establishes
the every-center upper bound. A selected positive squared norm is not
a lower bound for the complete squared norm: the complementary vector
can cancel it, and the complementary energy is still unknown.
